\begin{prob}{004}{\hosi 2}{}
 格子点を3頂点に持つ正三角形は存在しないことを示せ。また、「ピックの定理」を用いた証明も考えよ。
\end{prob}

複素数平面で考える. このとき格子点は実部虚部が整数の複素数(ガウス整数)に対応する. 3頂点がガウス整数の点のなる正三角形が存在すると仮定しよう. このとき, 適切に平行移動を施すことにより, ある頂点は複素数$0$であるとしてよい. 残りの頂点は整数$a,b,c,d$を用いて
$$
a+bi, c+di
$$
と表すことができ, $a+bi$ を0中心に反時計回りに$60^\circ$ 回転した点が$c+di$であるとしてよい. このとき, 複素数としては
$$
c+di = (\cos{60^\circ} + i\sin{60^\circ})(a+bi) = \dfrac{a-b\sqrt{3}}{2} + \dfrac{\sqrt{3}a+b}{2}i
$$
が成り立つから, 最右辺の実部と虚部が整数であるためには, 特に有理数でなければならないから, $\sqrt{3}$ が無理数であることを用いることにより $a=b=0$ が従う. これは, 異なる頂点がどちらも複素数0に対応していることになるため, 矛盾である. よって, このような正三角形は存在しないことが示された.

(\textbf{ピックの定理})

格子点を頂点に持つ多角形の面積$S$は,
内部の格子点の個数を$i$, 辺上の格子点の個数を$b$とするとき,
$S = i + \frac{b}{2} - 1$である.
特に, $S$は有理数である.

一方で, 格子点を頂点に持つ正三角形が存在するとしたとき, その1辺の長さ$a$は, 整数$m,n$ によって
$$
a = \sqrt{m^2 + n^2}
$$
と表すことができる. ただし, $(m,n)\neq (0,0)$である. このとき正三角形の面積は
$$
\frac{\sqrt{3}}{4} a^2 = \dfrac{\sqrt{3}}{4}(m^2 + n^2)
$$
であるから, $\sqrt{3}$ が無理数であることを用いることにより,
この面積は有理数でなく, 矛盾である. よって, 格子点を頂点に持つ正三角形は存在しないことが示された.
